% Бүлэг

\chapter{Удиртгал}
\label{ChapterIntro}

Дижитал хөгжмийн үйлдвэрлэл хөгжихийн хэрээр FL Studio зэрэг дижитал аудио ажлын орчин (DAW) ашиглан хэмнэл бүтээх, ая зохиох, бүтээлийн зохиомж боловсруулах, холилт болон мастеринг хийх сонирхолтой хэрэглэгчдийн тоо нэмэгдэж байна. Гэвч эхлэн суралцагчдын хувьд хөгжмийн онол, програмын хэрэгсэл, бүтээлч ажлын урсгал, техникийн сонголтуудыг нэгэн зэрэг ойлгох шаардлагатай байдаг тул суралцах үйл явц ихэвчлэн тасралттай, тодорхой дараалалгүй болдог. Интернэт орчинд хөгжмийн сургалтын материал олон боловч тэдгээр нь чанар, түвшин, дараалал, дадлагын чиглүүлэг, ахицын хяналтын хувьд харилцан адилгүй байна.

Ийм нөхцөлд суралцагч зөвхөн тусдаа видео үзэх бус, харин мэргэжлийн багшийн бэлтгэсэн хичээл, дараалсан дадлага, өөрийгөө шалгах сорил, төлбөртэй хандалтын удирдлага, хувийн ахицын хяналт, шаардлагатай үед асуултаа лавлах AI туслахыг нэг орчинд ашиглах хэрэгцээ үүсэж байна. Энэхүү дипломын ажил нь уг хэрэгцээнд тулгуурлан FL Studio болон хөгжим боловсруулалт сурах хэрэглэгчдэд зориулсан, хиймэл оюуны дэмжлэгтэй веб сургалтын платформын дизайн ба хөгжүүлэлтийг авч үзнэ.

\section*{Зорилго}
Энэхүү дипломын ажлын зорилго нь FL Studio болон хөгжим боловсруулалт суралцагчдад зориулан мэргэжлийн видео хичээл, төлбөртэй хичээлийн хандалт, хичээлийн тоглуулагч, өөрийгөө шалгах сорил, суралцах ахицын хяналт болон локал RAG технологид суурилсан сургалтын туслахыг нэгтгэсэн веб сургалтын платформ зохион бүтээж, хөгжүүлэхэд оршино.

\section*{Зорилт}
Энэхүү зорилгыг хэрэгжүүлэхийн тулд дараах зорилтуудыг дэвшүүлэв:
\begin{itemize}
    \item FL Studio болон хөгжим боловсруулалт суралцагчдад тулгардаг үндсэн хүндрэл, онлайн сургалтын хэрэгцээг тодорхойлох;
    \item ижил төстэй сургалтын платформ, AI туслах платформ болон хөгжмийн AI судалгааны ажлуудыг харьцуулан судлах;
    \item курсийн каталог, төлбөртэй хичээлийн хандалт, хичээлийн тоглуулагч, хяналтын самбар, өөрийгөө шалгах сорил, ахицын хяналт бүхий сургалтын урсгалыг загварчлах;
    \item Next.js, React, TypeScript, Supabase болон PostgreSQL-д тулгуурласан веб платформын архитектур, өгөгдлийн загварыг боловсруулах;
    \item FL Studio, хэмнэл бүтээх арга, холилт, мастеринг болон хөгжмийн онолын асуултад контекстэд тулгуурлан хариулах Ollama/Qdrant суурьтай RAG чатботыг платформд нэгтгэх;
    \item хөгжүүлсэн платформын үндсэн ажиллагаа, хэрэглэгчийн урсгал, өгөгдлийн хадгалалт болон AI туслахын хэрэглээний боломжийг турших.
\end{itemize}

\section*{Платформ хөгжүүлэх үндэслэл}
Хөгжим боловсруулалтын сургалтын платформ хөгжүүлэх хэрэгцээ нь сургалтын зохион байгуулалт, хэрэглэгчийн туршлага, технологийн боломж гэсэн гурван үндэслэлээр тайлбарлагдана.

Платформ хөгжүүлэх гол үндэслэл нь:
\begin{enumerate}
    \item FL Studio суралцах үйл явц нь хэмнэл бүтээх, аялгуу ба гармони боловсруулах, зохиомж хийх, холилт, мастеринг зэрэг олон чадварыг үе шаттай эзэмшихийг шаарддаг. Иймд тархай эх сурвалжаас суралцахаас илүү дараалсан сургалтын хөтөлбөр, дадлага, өөрийгөө шалгах сорил бүхий сургалтын замнал шаардлагатай.
    \item Багш, продюсерын бэлтгэсэн төлбөртэй видео хичээлийг хэрэглэгчийн эрхийн удирдлага, хичээлийн явцын бүртгэл, хяналтын самбар болон өгөгдлийн сангийн найдвартай бүтэцтэй холбосноор сургалтын чанар, хүртээмжийг нэмэгдүүлэх боломжтой.
    \item Орчин үеийн хэлний загвар болон мэдээлэл сэргээн баяжуулсан генерацийн аргачлал нь суралцагчийн асуултад гаднын мэдлэгийн сангаас сонгосон контекст ашиглан хариулах боломжийг бүрдүүлж байна \cite{attention_vaswani_2017,rag_lewis_2020,dpr_karpukhin_2020}. Энэ нь хичээл үзэх явцад гарсан тодорхой асуултыг шуурхай тайлбарлах нэмэлт сургалтын хэрэгсэл болж чадна.
\end{enumerate}

Иймээс энэхүү платформ нь хөгжмийн сургалтын агуулгыг бүтэцтэй хүргэх, суралцагчийн ахицыг хэмжих, төлбөртэй контентын хандалтыг зохион байгуулах, мөн AI туслахыг хязгаарлагдсан сургалтын контекстэд ашиглах бодит шаардлагад тулгуурлаж байна.

\section*{Судалгааны өнөөгийн байдал}
Сүүлийн жилүүдэд Transformer архитектур том хэлний загварын үндсэн суурь болж, урт context болон асуулт-хариултын даалгаварт өргөн ашиглагдах болсон \cite{attention_vaswani_2017}. Том хэлний загварууд few-shot болон instruction-following чадвараараа сургалтын туслах системд ашиглах боломжийг нэмэгдүүлсэн \cite{gpt3_brown_2020,llama_touvron_2023}.

Мэдээлэл сэргээн баяжуулсан генераци буюу RAG нь хэлний загварыг гаднын мэдлэгийн сантай холбож, хариултыг сонгосон context дээр тулгуурлуулах арга юм \cite{rag_lewis_2020}. Энэ хандлага нь боловсролын платформд буруу, ерөнхий эсвэл эх сурвалжгүй зөвлөгөө өгөх эрсдэлийг багасгахад тохиромжтой \cite{atlas_izacard_2022}.

Хөгжим-AI чиглэлд текстээс хөгжим үүсгэх, audio representation суралцах, яриа болон аудио боловсруулах судалгаанууд эрчимтэй хөгжиж байна \cite{jukebox_dhariwal_2020,audiolm_borsos_2022,musicgen_2023,musiclm_2023}. Мөн wav2vec 2.0 болон Whisper зэрэг ажлууд аудио-текстийн боловсруулалтад хүчтэй суурь болсон \cite{wav2vec2_baevski_2020,whisper_radford_2022}. Гэвч энэхүү дипломын ажлын хүрээнд AI-г бүрэн автомат хөгжим үүсгэгч бус, харин мэргэжлийн багшийн боловсруулсан сургалтын агуулгыг дэмжих, хичээлтэй холбоотой асуултад контексттэй хариулах туслах функц байдлаар авч үзсэн.

\section*{Шинэлэг тал}
Энэхүү платформын шинэлэг тал нь FL Studio болон хөгжим боловсруулалтын төлбөртэй видео сургалтыг локал RAG туслах, хэрэглэгчийн ахицын хяналт, хичээлийн эрхийн удирдлагатай нэг платформд уялдуулсанд оршино. Үүнд:
\begin{itemize}
    \item хичээл, дадлага, сорил, ахицын мэдээллийг нэг сургалтын урсгалд холбосон хэрэглэгчийн интерфейс;
    \item төлбөртэй хичээлийн хандалт, хэрэглэгчийн бүртгэл, хяналтын самбар болон ахицын хяналтыг нэг өгөгдлийн загварт нэгтгэсэн шийдэл;
    \item FL Studio болон хөгжмийн онолын мэдлэгийн сангаас хэрэглэгчийн асуултад тохирох контекст сонгон хариулах локал RAG туслах;
    \item мэргэжлийн багшийн хичээлд суурилсан агуулгыг эх сурвалжтай зөвлөмжөөр нэмэлтээр дэмжсэн сургалтын туршлага.
\end{itemize}

\section*{Ач холбогдол}
Энэхүү платформыг хэрэгжүүлснээр FL Studio болон хөгжим боловсруулалт сурах хүсэлтэй хэрэглэгчид сургалтын агуулгыг шат дараатай үзэх, өөрийн ойлголтоо сорилоор шалгах, хичээлийн явцаа хянах, тодорхой асуудал гарсан үед AI туслахтай харилцан лавлах боломжтой болно. Энэ нь эхлэн суралцагчдын суралцах тасралтыг бууруулж, дадлагад суурилсан чадвар эзэмшилтийг дэмжинэ.

Мөн багш, продюсеруудад мэргэжлийн контентоо бүтэцтэй сургалтын бүтээгдэхүүн болгон хүргэх, төлбөртэй хандалтаар орлогожуулах, суралцагчийн ахицын мэдээлэлд тулгуурлан сургалтын агуулгаа сайжруулах боломж бүрдэнэ. Бүтээлч салбарын боловсролын хувьд уг ажил нь веб технологи, өгөгдлийн сан, хэрэглэгчийн туршлага болон AI туслах боломжийг нэгтгэсэн жишээ загвар болох практик ач холбогдолтой.

\section*{Хамрах хүрээ}
Энэхүү дипломын ажлын хамрах хүрээнд хэрэглэгчийн бүртгэл, нэвтрэлт, курсийн жагсаалт, төлбөртэй хичээлийн хандалт, хичээлийн тоглуулагч, өөрийгөө шалгах сорил, суралцах явцын хяналт, хяналтын самбар, админ удирдлага, мөн хичээлтэй холбоотой асуултад контекстээр дэмжлэг үзүүлэх AI чатбот багтана. Платформ нь FL Studio болон хөгжим боловсруулалтын сургалтын хэрэглээнд төвлөрөх бөгөөд хүний багшийг бүрэн орлох, хэрэглэгчийн өмнөөс бүрэн автомат хөгжим бүтээх, эсвэл студийн бүх үйлдвэрлэлийн процессыг автоматжуулах зорилгогүй.

\section*{Тайлангийн бүтэц}
Энэхүү тайлангийн бүтэц нь судалгааны үндэслэлээс эхлэн технологийн сонголт, архитектур, хэрэгжилт, туршилтын үр дүн хүртэл үе шаттайгаар уялдах зарчимд тулгуурлав. Нэгдүгээр бүлэгт төслийн үндэслэл, ижил төстэй платформ, шаардлага болон хөгжүүлэх арга зүйг тайлбарлана. Хоёрдугаар бүлэгт веб технологи, өгөгдлийн сан, AI, Transformer, RAG болон хөгжмийн AI-ийн үндсийг судална. Гуравдугаар бүлэгт платформын архитектур, фронтенд, бекенд, өгөгдлийн сан, AI чатбот, өгөгдлийн багц бэлтгэгч зэрэг хэрэгжилтийн хэсгүүдийг нэгтгэн тайлбарлана. Дөрөвдүгээр бүлэгт туршилт, хэрэглэгчийн үндсэн урсгал, чатботын хариу, өгөгдлийн чанар болон үр дүнг авч үзнэ.
